\section{Introduction}
\label{sec:introduction}

\subsection{The Decadal Reapportionment Question}

The \dia{} specifies that the \ar{} algorithm is run fresh after each
decadal census, using the new \hh{} apportionment~\citep{huntington1928}
as the district count $k$ for each state.
This is the intended behavior: reapportionment is a constitutionally mandated
response to population change~\citep{wesberry1964}, and the algorithm
should reflect the new apportionment.

But what actually happens when a state gains or loses seats?
The \ar{} bisection tree is determined entirely by the prime factorisation
of $k$ (B.11~\citep{b11paper}).
A change in $k$ can dramatically alter the tree topology---not because the
algorithm has changed, but because the arithmetic has.

This paper characterises the 2030 reapportionment shock.
We use Census Bureau 2030 projections~\citep{census2030proj} to identify
which states will gain or lose seats, trace how each seat change alters
the \ar{} bisection tree, and simulate the resulting tract-level boundary
movements.

\subsection{The Key Cases}

\paragraph{Texas ($38 \to 41$).}
$38 = 2 \times 19$ is a two-level tree (bisect to 2, each half bisects
to 19 by recursive bisection since 19 is prime).
$41$ is prime: the entire tree is a flat 41-way split at the root.
This is the most dramatic factorisation change in the nation.
We expect large boundary movement.

\paragraph{California ($52 \to 50$).}
$52 = 4 \times 13 = 2^2 \times 13$.
$50 = 2 \times 25 = 2 \times 5^2$.
Both are composite, but with different tree structures.
We characterise the structural change and expected boundary movement.

\paragraph{Florida ($28 \to 29$).}
$28 = 4 \times 7 = 2^2 \times 7$.
$29$ is prime: flat tree, major structural change.

\paragraph{New York ($26 \to 25$).}
$26 = 2 \times 13$.
$25 = 5^2 = 5 \times 5$: balanced two-level tree.

\subsection{Main Finding}

The main finding: seat-change-induced boundary disruption is \emph{less}
than the disruption from a single \recom{} MCMC remix step.
This makes intuitive sense: the \ar{} algorithm is optimising for compactness
within the new apportionment, which requires finding a new optimal map
from scratch, but the geographic constraints are largely unchanged.

\subsection{Paper Structure}

Section~\ref{sec:projections} summarises 2030 reapportionment projections.
Section~\ref{sec:factorization} analyses the prime factorisation changes.
Section~\ref{sec:boundary} simulates boundary stability.
Section~\ref{sec:gerrychain-baseline} establishes the GerryChain baseline.
Section~\ref{sec:statutory} discusses the statutory implications.
Section~\ref{sec:conclusion} concludes.
